\begin{table}[!h]
\centering
\caption{\label{tab:p8_tcn}Respuesta de expectativas de tipo de cambio al shock cambiario}
\centering
\begin{threeparttable}
\begin{tabular}[t]{cccc}
\toprule
$h$ & Expectativa & Realizado 12m fwd & Gap (exp $-$ real)\\
\midrule
0 & -0.830 [-0.927, -0.732] & -0.084 [-0.494, 0.326] & -0.669 [-1.065, -0.274]\\
3 & -0.226 [-0.383, -0.069] & -0.208 [-0.660, 0.244] & 0.018 [-0.380, 0.415]\\
6 & -0.090 [-0.206, 0.027] & -0.282 [-0.694, 0.129] & 0.226 [-0.181, 0.634]\\
12 & 0.032 [-0.145, 0.209] & 0.118 [-0.241, 0.478] & -0.096 [-0.500, 0.309]\\
18 & 0.046 [-0.103, 0.195] & 0.053 [-0.334, 0.441] & -0.004 [-0.377, 0.368]\\
\addlinespace
24 & 0.058 [-0.052, 0.168] & -0.044 [-0.449, 0.362] & 0.109 [-0.299, 0.517]\\
\bottomrule
\end{tabular}
\begin{tablenotes}
\item \textit{Note: } 
\item Coeficiente $beta_h$ de la LP con el shock cambiario $hat u_t$ del VAR bivariado (Secci'on 3) como regresor; controles: rezagos del VAR. emph{Expectativa}: respuesta de la depreciaci'on esperada a 12 meses, $E_{t+h}[tcn_{t+h+12}]-tcn_{t+h}$, construida como $text{exp_cam}_{t+h}-tcn_{t+h}$. emph{Realizado 12m fwd}: respuesta de $tcn_{t+h+12}-tcn_{t+h}$, la depreciaci'on efectivamente realizada, comparable con la expectativa. emph{Gap}: diferencia entre ambas; su coeficiente es el test directo de igualdad. Intervalos al 95% con errores de Newey-West y $h+1$ rezagos. Ambas series en puntos log a 12 meses. Muestra: 2001m11--2025m12.
\end{tablenotes}
\end{threeparttable}
\end{table}